%Revisado
\subsection{Limitações Sensoriais e Multiespectrais}

    Embora câmeras RGB sejam predominantes devido ao baixo custo, a literatura aponta limitações físicas críticas na detecção de certos contaminantes. \citeonline{mustafic2016} demonstraram que materiais como plásticos transparentes e papéis brancos possuem assinaturas visuais quase idênticas à da fibra de algodão no espectro visível, tornando-os virtualmente invisíveis para sensores convencionais. A solução proposta pelos autores envolve o uso de imageamento hiperespectral de fluorescência (UV), capaz de distinguir a composição química de contaminantes não botânicos com precisão superior a 90\%.

    Paralelamente, \citeonline{li2024cottonnet} exploraram o uso de espectroscopia no infravermelho próximo (NIR) combinada com redes neurais (\textit{Cotton-Net}) para estimar o teor de impurezas. Embora o método seja eficaz para determinar a composição química global da amostra, ele carece de informação espacial detalhada para classificar a morfologia dos defeitos, reafirmando a necessidade da visão computacional para a avaliação física do aspecto visual do fardo.

\subsection{O Aprendizado de Máquina e a Variabilidade Intra-Amostra}

    A transição da estatística clássica para o aprendizado de máquina marcou a primeira etapa da automação robusta. \citeonline{fisher2023} investigaram a classificação de algodão egípcio utilizando algoritmos de \textit{Random Forest}, superando abordagens baseadas em Máquinas de Vetores de Suporte (\textit{Support Vector Machines}) e redes neurais artificiais clássicas. Uma contribuição fundamental desse estudo foi a introdução da variância intra-amostra como preditor de qualidade: os autores provaram que a uniformidade da cor ao longo da superfície da amostra é tão determinante para a classificação comercial quanto a média dos valores de reflectância.

    Contudo, a eficácia desses modelos esbarra em um gargalo humano: a qualidade da rotulagem dos dados (\textit{ground truth}). \citeonline{fisher2023aiassisted} alertam que a subjetividade e a fadiga dos classificadores humanos introduzem ruído nos dados de treinamento. Para mitigar o alto custo e a inconsistência da anotação manual, os autores propuseram o uso de aprendizado ativo (\textit{Active Learning}), no qual o sistema solicita intervenção humana apenas para imagens de alta incerteza, reduzindo drasticamente o volume de dados necessários enquanto mantém a acurácia do modelo.

\subsection{Abordagens Contemporâneas de Visão Computacional para o Algodão}

    Atualmente, a fronteira do conhecimento reside na aplicação de CNNs para tarefas de segmentação semântica e detecção de objetos, com foco na correlação direta com os parâmetros do HVI.

    No que tange à análise de cor, o principal desafio é a interferência de sombras e sujeira na leitura dos pixels. \citeonline{li2025measurement} propuseram o método inovador \textit{Color-Unet}. Diferentemente das abordagens que calculam a média de toda a imagem, essa rede realiza a segmentação semântica para isolar a fibra pura, excluindo pixels de impurezas e áreas sombreadas. A medição de cor realizada apenas na máscara de fibra limpa resultou em uma correlação significativamente maior com os valores de referência do HVI.

    Para a detecção de impurezas (\textit{Trash}), a arquitetura YOLO (\textit{You Only Look Once}) estabeleceu-se como padrão. \citeonline{li2025cottonyolo} desenvolveram o modelo \textit{Cotton-YOLO}, uma versão otimizada com módulos de atenção projetados especificamente para detectar impurezas milimétricas que redes genéricas não conseguem identificar. Avançando para a quantificação precisa, \citeonline{jiang2024cottonyoloseg} introduziram o \textit{Cotton-YOLO-Seg}, que não apenas detecta o objeto, mas realiza a segmentação de instância (pixel a pixel). Isso permite o cálculo exato da área ocupada por cada partícula, viabilizando a automação do parâmetro \textit{\% Area} do HVI com precisão inédita.

    Por fim, a avaliação da preparação (irregularidades como nós de fibra) também se beneficia dessas arquiteturas. \citeonline{arachchi2024classification} aplicaram CNNs baseadas em \textit{transfer learning} (como a \textit{Inception V3}) para classificar defeitos em fios de algodão, demonstrando que padrões morfológicos sutis de emaranhamento podem ser aprendidos eficazmente por redes profundas, uma lógica que se estende à avaliação da pluma crua.